\section{Completion and the generic specialization kernel}
\label{sec:completion-generic-kernel}

Completion need not algebraize every coefficient-free representation.
It is injective on the entire lattice Grothendieck group and preserves
the entire generic specialization kernel.
The distinction is useful here: after adding one relative syzygy,
every class in that kernel can be represented by lattices whose
generic group-algebra modules are free, and those lattices admit a
direct algebraization by approximation of their coefficient bases.

\begin{theorem}\label{thm:completion-generic-kernel}\label{thm:completion-dvr-kernel}
Let $O$ be any characteristic-zero DVR, with uniformizer $\pi$,
fraction field $K$, completion $\widehat O$, and completed fraction
field $\widehat K$.  Let $A$ be a nonzero commutative algebra finite
free over $O$, and let $G$ be a finite group.  Put
\[
 \widehat A=A\otimes_O\widehat O,\qquad
 D=A[1/\pi],\qquad \widehat D=\widehat A[1/\pi].
\]
Completion induces an injection
\[
 K_0(\Lat_A(G))\hookrightarrow K_0(\Lat_{\widehat A}(G))
\]
and an isomorphism
\[
 \ker\!\left(K_0(\Lat_A(G))\longrightarrow K_0(\Lat_D(G))\right)
 \xrightarrow{\ \sim\ }
 \ker\!\left(K_0(\Lat_{\widehat A}(G))
                   \longrightarrow K_0(\Lat_{\widehat D}(G))\right).
\]
Moreover, any subgroup of the right-hand kernel that is split in
$K_0(\Lat_{\widehat A}(G))$ lifts to a subgroup of the left-hand
kernel that is split in $K_0(\Lat_A(G))$.
\end{theorem}

Here $\Lat_A(G)$ consists of finite free $A$-modules with
$A$-linear $G$-action, and its admissible sequences are the
coefficient-split short exact sequences.  A free module has one common
finite rank over the whole coefficient ring, including when that ring
is nonlocal or has several components.  Neither locality nor
reducedness of $A$ is assumed, and no Henselian, excellence,
ramification, or residue-field condition is imposed on $O$.  The theorem makes no assertion that all
completed lattices descend, nor that completion is an equivalence
on morphism groups.

\subsection{Generic projectives and algebraization}

We use first an elementary observation about the generic categories.
If $C$ is a finite-dimensional commutative algebra over a
characteristic-zero field, every $C$-free $G$-module is projective
over $C[G]$: a coefficient-linear section of an epimorphism onto it
becomes equivariant after averaging over $G$.  Thus every admissible
sequence in $\Lat_C(G)$ splits equivariantly.  The algebra $C[G]$
is Artinian, and its finitely generated projective modules have
unique multiplicity vectors in the finite set of indecomposable
projective classes.  The coefficient-free projectives form a
submonoid of this free commutative monoid.  Its group completion
embeds in the corresponding free abelian group.  Consequently
\begin{equation}\label{eq:completion-generic-cancellation}
 K_0(\Lat_C(G))\text{ is torsion-free},\qquad
 [U]=[V]\text{ in }K_0(\Lat_C(G))\Longrightarrow U\cong V.
\end{equation}
This includes nonreduced coefficient algebras $C$.

\begin{lemma}\label{lem:completion-generic-free-algebraization}
Let $M\in\Lat_{\widehat A}(G)$.  If
$M[1/\pi]\cong\widehat D[G]^s$, then there exists
$L\in\Lat_A(G)$ with
\[
 \widehat A\otimes_A L\cong M,
 \qquad L[1/\pi]\cong D[G]^s.
\]
\end{lemma}

\begin{proof}
Choose the displayed generic isomorphism and view $M$ inside the
standard representation $\widehat D[G]^s$.  With $k=s|G|$, a
coefficient basis of $M$ gives
\[
 M=P\widehat A^k,\qquad P\in\mathrm{GL}_k(\widehat D).
\]
The algebra $D$ is dense in $\widehat D$: in an $O$-basis of
$A$ this is the density of $K$ in $\widehat K$.  Choose
$P_0\in\mathrm M_k(D)$ sufficiently close to $P$ that
\[
 P^{-1}P_0\in I+\pi\mathrm M_k(\widehat A).
\]
Then $P_0\widehat A^k=M$.  The matrix $P_0$ is invertible over
$D$, since it is invertible after the faithfully flat scalar
extension $K\to\widehat K$.

Set $L=P_0A^k\subset D[G]^s$.  Let $R_g$ denote the standard
regular-action matrix of $g\in G$.  Stability of $M$ gives
\[
 P_0^{-1}R_gP_0\in\mathrm M_k(D)\cap
                         \mathrm M_k(\widehat A)
                       =\mathrm M_k(A).
\]
The last equality follows coordinatewise from
$K\cap\widehat O=O$ in any $O$-basis of $A$: membership in
the completed DVR is nonnegative valuation, and the valuation on $K$
is the original DVR valuation.
Applying the same statement to $g^{-1}$ shows that these matrices
are invertible over $A$.  Thus $L$ is a coefficient-free
$G$-lattice, its completion is $M$, and its generic module is the
chosen standard $D[G]^s$.
\end{proof}

\subsection{Reflection of isomorphisms and exact triples}

The completion facts used below are finite-module exactness,
flatness, faithful flatness when the completed ideal lies in the
Jacobson radical, and equality of finite quotients
\citep[Lemmas~10.97.1--10.97.4, Tags 00MA, 00MB, 00MC, 031C]{Stacks2026}.

\begin{lemma}\label{lem:completion-reflects-triples}
For coefficient-free $A$-lattices, completion reflects
isomorphisms and the existence of coefficient-split short exact
sequences with three prescribed isomorphism classes.
\end{lemma}

\begin{proof}
For $L,N\in\Lat_A(G)$, equivariant homomorphisms are the kernel
of the finitely many linear equivariance equations on
$\operatorname{Hom}_A(L,N)$.  Flatness of $\widehat O$ over $O$
therefore gives
\begin{equation}\label{eq:completion-hom-base-change}
 \operatorname{Hom}_{A[G]}(L,N)\otimes_O\widehat O
 \cong
 \operatorname{Hom}_{\widehat A[G]}(\widehat L,\widehat N).
\end{equation}
The integral Hom module is finite and torsion-free over the DVR
$O$, hence dense in its completion.  If $\widehat L\cong\widehat N$,
approximate an equivariant completed isomorphism modulo $\pi$ by an
integral equivariant map $U$.  Its reduction modulo $\pi$ is
invertible.  Since $A$ is finite over $O$, the ideal $\pi A$ lies in
its Jacobson radical, and $A/\pi A=\widehat A/\pi\widehat A$.
Therefore $\det U$ is a unit in $A$, so $U$ is an isomorphism.

For the exact-sequence assertion, coefficient-split extensions of
$N$ by $L$ are classified by
\[
 H^1\!\left(G,\operatorname{Hom}_A(N,L)\right),
\]
where $G$ acts by conjugation.  Indeed, a coefficient-linear
splitting writes the action in upper block-triangular form, whose
off-diagonal entries give a one-cocycle; changing the splitting
adds a coboundary.

This cohomology group is finite over $O$ and is killed by $|G|$.
For an explicit check of the latter statement, if $c_g$ is a
one-cocycle and $S=\sum_hc_h$, summing
$c_{gh}=c_g+gc_h$ over $h$ yields
\[
 |G|c_g=S-gS.
\]
Since $O$ has characteristic zero, $|G|$ is nonzero in $O$; it is
a unit times a power of $\pi$.  Thus this cohomology group has finite
length over $O$ (and is zero if $|G|$ is a unit), so it is unchanged
by completion.  Flat base change of the finite group cochain complex
also gives
\[
 H^1\!\left(G,\operatorname{Hom}_A(N,L)\right)\otimes_O\widehat O
 \cong
 H^1\!\left(G,
       \operatorname{Hom}_{\widehat A}(\widehat N,\widehat L)
       \right).
\]
Every completed extension class consequently descends to a
coefficient-split extension over $A$.  If its completed middle
term is isomorphic to the completion of a prescribed lattice,
the isomorphism-reflection argument just proved identifies the
integral middle term with that lattice.  This proves reflection
of exact triples.  Preservation follows from flat base change.
\end{proof}

\subsection{Cofinality and exact-sequence relations}

We use the following published cofinality lemma
\citep[Chapter~II, Lemma~7.2, p.~62 of the author's chapter
version]{Weibel2013}.

\begin{lemma}[Cofinality]\label{lem:completion-elementary-cofinality}
Let $\mathcal F$ be a full, isomorphism-closed, extension-closed
exact subcategory of an exact category $\mathcal E$, with the
induced exact structure.  Suppose every $X\in\mathcal E$ has a
supplement $X'\in\mathcal E$ such that
$X\oplus X'\in\mathcal F$.  Then
$K_0(\mathcal F)\to K_0(\mathcal E)$ is injective.
\end{lemma}

The extension-closure hypothesis controls exact-sequence relations
as well as stable isomorphisms.

\subsection{Identification of the kernels}

\begin{proof}[Proof of Theorem~\ref{thm:completion-generic-kernel}]
Write $\mathcal E_A=\Lat_A(G)$ and let $\mathcal F_A$ be its
full subcategory of lattices $L$ for which $L[1/\pi]$ is free over
$D[G]$.  Define $\mathcal E_{\widehat A}$ and
$\mathcal F_{\widehat A}$ in the same way.  Both subcategories
$\mathcal F$ are extension-closed: a generic extension of free
group-algebra modules splits and is free.

Lemmas~\ref{lem:completion-generic-free-algebraization}
and~\ref{lem:completion-reflects-triples} show that completion
gives a bijection between the isomorphism classes in
$\mathcal F_A$ and $\mathcal F_{\widehat A}$ and preserves and
reflects their exact triples.  Comparing the generators and
relations defining their Grothendieck groups gives
\begin{equation}\label{eq:completion-free-subcategory-kzero}
 K_0(\mathcal F_A)\xrightarrow{\sim}
 K_0(\mathcal F_{\widehat A}).
\end{equation}

These are cofinal subcategories in the sense of
Lemma~\ref{lem:completion-elementary-cofinality}.  Indeed, for
$M\in\mathcal E_A$ of coefficient rank $a$, the action map
\[
 A[G]\otimes_A M\longrightarrow M,
 \qquad g\otimes m\longmapsto gm,
\]
is a coefficient-split epimorphism.  Its kernel $\Omega M$ is
coefficient-free: an explicit coefficient isomorphism from
$\bigoplus_{g\ne1}M$ onto that kernel is
\[
 (m_g)_{g\ne1}\longmapsto
 \sum_{g\ne1}\bigl(g\otimes m_g-1\otimes gm_g\bigr).
\]
The middle term is $A[G]^a$.  Generically $M[1/\pi]$ is projective
by averaging, so
\[
 M[1/\pi]\oplus(\Omega M)[1/\pi]\cong D[G]^a.
\]
Consequently $M\oplus\Omega M\in\mathcal F_A$.  The same
argument applies after completion.  The lemma now gives
injections
\begin{equation}\label{eq:completion-cofinal-injections}
 K_0(\mathcal F_A)\hookrightarrow K_0(\mathcal E_A),\qquad
 K_0(\mathcal F_{\widehat A})\hookrightarrow
                         K_0(\mathcal E_{\widehat A}).
\end{equation}

We can now prove injectivity on the full Grothendieck groups.
Let $\mathcal I_A$ be the full, isomorphism-closed subcategory of
$\mathcal E_{\widehat A}$ consisting of objects isomorphic to
completions of objects of $\mathcal E_A$.  The extension-descent
argument in Lemma~\ref{lem:completion-reflects-triples} shows that
$\mathcal I_A$ is extension-closed.  Reflection of isomorphisms
and preservation and reflection of exact triples identify the
generators and relations, giving
\[
 K_0(\mathcal E_A)\xrightarrow{\sim}K_0(\mathcal I_A).
\]
For every $M\in\mathcal E_{\widehat A}$, the explicit syzygy above
gives $M\oplus\Omega M\in\mathcal F_{\widehat A}$.
Lemma~\ref{lem:completion-generic-free-algebraization} gives
$\mathcal F_{\widehat A}\subseteq\mathcal I_A$.  Thus
$\mathcal I_A$ is cofinal in $\mathcal E_{\widehat A}$, and
Lemma~\ref{lem:completion-elementary-cofinality} yields
\[
 K_0(\mathcal E_A)\xrightarrow{\sim}K_0(\mathcal I_A)
 \hookrightarrow K_0(\mathcal E_{\widehat A}).
\]
This composite is completion.  The argument uses algebraization
only for regular generic type and makes no descent assertion for
an arbitrary completed lattice.

The generic free-group-algebra rank defines homomorphisms
\[
 \nu_A:K_0(\mathcal F_A)\to\Z,
 \qquad
 \nu_{\widehat A}:K_0(\mathcal F_{\widehat A})\to\Z,
\]
and the isomorphism
\eqref{eq:completion-free-subcategory-kzero} preserves them.
Under \eqref{eq:completion-cofinal-injections},
$\ker\nu_A$ is exactly the generic specialization kernel in
$K_0(\mathcal E_A)$.  To see surjectivity onto that kernel,
write any class as $[M]-[N]$.  If its generic image is zero,
\eqref{eq:completion-generic-cancellation} gives
$M[1/\pi]\cong N[1/\pi]$.  Adding $\Omega M$ to both terms
represents the same class by
\[
 [M\oplus\Omega M]-[N\oplus\Omega M],
\]
and both lattices are in $\mathcal F_A$ with the same generic
free-group-algebra rank.  Conversely, a difference of rank zero
in $K_0(\mathcal F_A)$ has zero generic class.  The same reasoning
identifies the completed kernel with $\ker\nu_{\widehat A}$.
Equation~\eqref{eq:completion-free-subcategory-kzero} therefore
proves the asserted isomorphism of kernels.

Finally, let $s:H\to K_0(\mathcal E_{\widehat A})$ and
$q:K_0(\mathcal E_{\widehat A})\to H$ satisfy $qs=1$, with
$s(H)$ contained in the completed generic kernel.  Lift $s$
through the kernel isomorphism, obtaining
$s_A:H\to K_0(\mathcal E_A)$, and put
\[
 q_A=q\circ\text{completion}.
\]
Then $q_As_A=1$.  The lifted subgroup is consequently split in
the entire integral lattice Grothendieck group, as claimed.
\end{proof}

\subsection{Deep monic deformations over the localized integers}

\begin{corollary}\label{cor:deep-monogenic-orders}
Let $O=\Z_p$ or $\Z_{(p)}$, and let $F\in O[t]$ be monic of
degree $n\ge2$ with
\[
 F\equiv t^n\pmod {p^3}.
\]
For $B=O[t]/F$, the generic specialization kernel in
$K_0(\Lat_B(S_p))$ contains a subgroup isomorphic to
$\bigoplus_{i\ge1}\Z/n$ that is split in the whole Grothendieck
group.  A countable such summand has a coordinate retraction
annihilating the Young permutation span.  In particular,
$\Phi_B$ is not surjective in degree $p$.
\end{corollary}

\begin{proof}
First let $O=\Z_p$, and put $B_0=\Z_p[t]/t^n$.
The ring maps $\Z_p\to B_0\to\Z_p$, with $t\mapsto0$ in the
second map, split the reduction homomorphism on lattice
Grothendieck groups.  Its kernel is consequently a direct summand
of the whole group.  The split countable $\Z/n$ subgroup of
Theorem~\ref{thm:all-truncations} is therefore split in
$K_0(\Lat_{B_0}(S_p))$.

The coefficient algebra $B$ is finite free of rank $n$, and the
hypothesis on $F$ gives $B/p^3B\cong B_0/p^3B_0$.
Since $v_p(|S_p|)=1$, Theorem~\ref{thm:congruent-coefficients}
applies with $r=3>2$.  Transport the preceding subgroup,
together with its section and retraction.  Its exact orders
and its splitting are preserved.  The generic target is torsion-free
by \eqref{eq:completion-generic-cancellation}, applied to
$B[1/p]$, so this torsion subgroup lies in the generic kernel.

For completeness, the Young-span assertion does not require a
particular normalization of the original retraction.  There are
only finitely many degree-$p$ Young permutation lattices.  Each
has finitely supported image in $\bigoplus_i\Z/n$.  Delete the
finite union of those supports and restrict the section to the
remaining coordinates.  Composing the retraction with their
coordinate projection gives a retraction of a countable direct
summand that annihilates every Young permutation lattice.
If the retraction from Theorem~\ref{thm:all-truncations} is already
so normalized, no coordinates need be deleted.

Now let $O=\Z_{(p)}$.  The completion of $B$ is
$\widehat B=\Z_p[t]/F$, so the complete case supplies a split
countable subgroup in its generic kernel.
Theorem~\ref{thm:completion-generic-kernel} lifts that subgroup
to the generic kernel over $B$ and preserves its splitting in
the full lattice Grothendieck group.  Compose the completed
coordinate retraction with completion.  Each integral Young
permutation lattice completes to its corresponding completed
Young lattice, so this retraction still annihilates the Young
span.  It takes the lifted subgroup onto the nonzero coordinate
vectors.  Since the degree-$p$ image of $\Phi_B$ is the Young
permutation span, $\Phi_B$ cannot be surjective.
\end{proof}
